% =========================================================
% Capstone Exercise: Rational Numbers
% =========================================================
% Status: Planned
% Dependency ceiling: all environments in rationals/chapter.yaml
% =========================================================

\newpage
\phantomsection
\label{cap:rationals}

\begin{tcolorbox}[
  colback=gray!6,
  colframe=gray!40,
  arc=2pt,
  left=8pt, right=8pt, top=6pt, bottom=6pt,
  title={\small\textbf{Capstone — Rational Numbers}},
  fonttitle=\small\bfseries
]
\textbf{Status: Planned.}

This capstone synthesizes the core results of the rationals chapter.  It is
designed to require the student to connect the algebraic, order-theoretic,
density, and analytic threads of the chapter in a single argument.

\medskip

\textbf{Problem (Draft).}
Let $(x_n)$ and $(y_n)$ be Cauchy sequences of rational numbers.

\begin{enumerate}
  \item Prove that $(x_n + y_n)$ and $(x_n y_n)$ are both Cauchy sequences.
  \item Define the relation $(x_n) \sim (y_n)$ if and only if
        $(x_n - y_n)$ converges to $0$ in $\mathbb{Q}$.
        Prove that $\sim$ is an equivalence relation on $\mathbb{Q}^c$.
  \item Suppose $(x_n) \sim (u_n)$ and $(y_n) \sim (v_n)$.
        Prove that $(x_n + y_n) \sim (u_n + v_n)$.
  \item Let $r \in \mathbb{Q}$ and let $(r)$ denote the constant sequence
        with every term equal to $r$.  Show that $(r) \in \mathbb{Q}^c$
        and that if $(x_n) \sim (r)$ then $\lim_{n \to \infty} x_n = r$.
  \item Construct an explicit Cauchy sequence $(x_n)$ in $\mathbb{Q}$ such
        that $(x_n) \not\sim (r)$ for any $r \in \mathbb{Q}$.  Identify
        which theorem in the chapter guarantees such a sequence exists.
\end{enumerate}

\medskip
\textit{Parts~1--3 are proved in the notes; this exercise requires the student
to reproduce and connect them in a single self-contained argument.
Part~4 makes precise the embedding $\mathbb{Q} \hookrightarrow \mathbb{Q}^c/{\sim}$.
Part~5 demonstrates incompleteness by explicit example.}

\end{tcolorbox}

\begin{remark*}[Dependency ceiling]
This capstone may use any result from the rationals chapter.
It may not invoke completeness of $\mathbb{R}$, Dedekind cuts,
or any result whose proof requires the real numbers.
\end{remark*}
